\label{chap:introduction}

Seabed detection is an intermediate step in echosounder processing, but it fixes a boundary that later analyses often treat as known. Fisheries acoustic surveys estimate biomass from water-column backscatter, and bottom recognition prevents the seabed echo from entering the integration \citep{thomas_trawl_2024,ona_mitson_1996}. A bottom line placed too shallow can discard fish near the seabed, whereas a line placed too deep can include seabed echo in the estimate. The detected bottom may also be the only operational depth reference available during a survey. In Norwegian territorial waters, public bathymetry is commonly limited to coarse 50\,m grids because finer data are restricted or only locally declassified \cite{kartverket_bathymetry_2024,mareano_bathymetric_2024}.

Marine echosounders obtain a depth estimate by emitting short acoustic pulses, or \emph{pings}, from a hull-mounted transducer. If the two-way travel time is $t$ and the speed of sound is $c$, the corresponding acoustic range is $d=ct/2$ (Figure~\ref{fig:echosounder_principle}). Each ping produces a vertical backscatter profile. Placing successive profiles side by side forms an \emph{echogram}, a two-dimensional range-ping image in which the seabed usually appears as a bright, mostly continuous band beneath the water column.

\begin{figure}[htbp]
    \centering
    \includegraphics[width=0.82\linewidth]{Figures/echosounder_diagram.png}
    \caption[Echosounder depth measurement.]{Echosounder depth measurement. A hull-mounted transducer transmits a pulse downward. The pulse reflects off the seabed and returns, and the two-way travel time gives the depth. As the vessel advances, successive measurements trace the seabed profile shown on the echogram display.}
    \label{fig:echosounder_principle}
\end{figure}

Traditional bottom detectors locate this band through threshold-based or energy-tracking rules. They usually process each ping as a one-dimensional echo profile and reduce features of the return to one depth estimate \citep{bourguignon_methodological_2009}. Such methods are fast, interpretable, and embedded in existing survey workflows. Their assumptions can nevertheless fail over steep or irregular terrain, near dense biological scattering, or when bubble interference, secondary echoes, and weak or distorted returns alter the expected bottom signature. An erroneous depth can then be returned without an internal indication that the estimate is unreliable.

An operator can often recognise the same failure by inspecting the echogram. Visual interpretation is not pointwise: continuity, local texture, and neighbouring pings help distinguish the seabed from isolated artefacts or biological returns. The mismatch between pingwise acoustic detection and image-based interpretation motivates a detector that reads two-dimensional echogram patches. Convolutional networks have performed well on segmentation tasks in medical imaging and are widely used in remote-sensing image analysis \citep{litjens_survey_2017,ma_deep_2019}. The echograms used here contain spatial dependence on a scale that a convolutional receptive field can exploit, as measured in Appendix~\ref{app:spatial-correlation}. Continuous echogram storage is already part of routine surveying, so model development can use data collected during normal operation rather than requiring a separate survey programme.

Deep learning has been applied to marine acoustics since at least 2017, including fish and species classification, subsea infrastructure detection, bottom detection in multibeam echosounder data, and bottom tracking in side-scan sonar imagery \citep{hirama_discriminating_2017,brautaset_acoustic_2020,yassir_acoustic_2023,du_revealing_2023,meng_antiinterference_2024,yan_bottom_2021}.

The proposed detector uses a U-Net \citep{ronneberger_unet_2015} to segment dual-frequency wideband Simrad EK80 echograms. Its input consists of calibrated backscatter at 38\,kHz and 200\,kHz together with the difference between the two channels. The network predicts a binary mask that separates the water column from the bottom-or-below region. Extracting the upper boundary of that mask gives one bottom-depth estimate per ping. Unlike a pingwise detector, this method can use neighbouring returns and the continuity of the seabed across the full image patch.

The primary data come from the M/F Bast\o{} Electric, which crosses an approximately 10.5\,km route in the Oslofjord about 24 times per day. The recordings were collected under the Frisk Oslofjord project \cite{frisk_oslofjord_mapping}. Repeated passes observe the same seabed under changing weather, sea state, tidal, biological, and acoustic conditions. Recording dates are separated between training, validation, conformal calibration, and testing, so the main held-out evaluation concerns later recordings from a known route rather than unseen geography. Repeated crossings also allow estimates from different passages to be compared at the same route location without requiring a manually drawn bottom line. The NOAA Pacific hake survey SH2306 provides a limited external transfer experiment in deeper water.

This design gives dense temporal coverage but restricts the claim that geographic transfer can be made. Performance on later Bast\o{} recordings tests variation within one operating area, one vessel installation, and one repeated bottom profile. It does not establish that an unchanged model transfers to other vessels, depth ranges, transducer configurations, or seabed types. The external survey is therefore treated as a stress test of transfer and recalibration, not as evidence for a universal detector.

Reference labels impose a second limitation. No independent seabed-depth measurements are available for the recordings, and systematic expert annotation was infeasible at the scale of the dataset and scope of the project. Training and supervised evaluation instead use masks generated from 38\,kHz acoustic returns by averaging a simple threshold-and-back-step detector with an edge detector and smoothing the resulting line over time \cite{crimac_classifiers_bottom}. These labels are algorithmic acoustic estimates, not measured seabed depths and not outputs from the operational EK80 detector. They can inherit errors when bubble layers, biological scatter, secondary echoes, or steep terrain obscure the bottom return. Supervised segmentation and depth scores therefore measure agreement with the reference pipeline. Likewise, conformal coverage against these labels quantifies model-reference discrepancy rather than error relative to the physical seabed. High agreement alone cannot show that either method has identified the true bottom.

The study addresses two questions. First, can a U-Net reproduce algorithmic seabed labels on later Bast\o{} recordings from the same route? Second, can the extracted depths be accompanied by calibrated intervals and diagnostics that identify unstable estimates? The evaluation combines segmentation overlap, depth agreement, probability diagnostics, supervised conformal intervals, repeated-crossing consistency, and external transfer. Supervised conformal intervals are calibrated from deviations between model and reference depths. The width of the probability transition around the predicted boundary is tested as a local scale for an adaptive interval, while the repeated-crossing analysis examines whether route-level disagreement can support a separate, label-free uncertainty band. The sigmoid output is therefore treated as an intermediate model output rather than as an uncertainty guarantee. The intended role is to supplement an existing detector with an alternative bottom estimate and reported uncertainty, not to replace operational EK80 detection. The model is therefore referred to throughout as the \emph{learned seabed detector}.

% MANDATORY SECTION: NTNU requires every thesis to relate its work to the UN
% Sustainable Development Goals. Do not remove this section.
\section{Relation to the UN Sustainable Development Goals}

The thesis relates to Goal 14, life below water, and target 14.4 on science-based management of fish stocks and the reduction of overfishing. Fisheries acoustic surveys provide data used in stock assessments, where the detected bottom defines the lower boundary of the integrated water column. A line placed above the seabed can exclude near-bottom fish, while a line placed below it can include seabed echo. Reliable detection with reported uncertainty can therefore reduce one source of error in the abundance estimates used to set fishing quotas.
